\section{Uso LLMs}

\noindent Para el desarrollo de este trabajo se utilizó una herramienta de LLM/IA como apoyo en tareas de revisión y mejora del texto.

\begin{enumerate}
    \item \textbf{Herramienta utilizada:} ChatGPT (GPT-5.5 Thinking).

    \item \textbf{Propósito de uso:} La herramienta fue utilizada principalmente para corregir faltas ortográficas, mejorar la redacción, revisar gramática, ordenar ideas y sugerir una formulación más clara de algunos párrafos del informe.

    \item \textbf{Prompts relevantes:} Se realizaron consultas orientadas a mejorar la calidad del texto, tales como: ``corrige la ortografía y redacción de este texto'', ``mejora la claridad de este párrafo'' y ``redacta esta sección en tercera persona y con estilo académico''.

    \item \textbf{Nivel de edición:} Las respuestas generadas fueron revisadas, adaptadas y modificadas antes de ser incorporadas al informe. No se copiaron sin revisión, ya que el contenido fue ajustado según los requerimientos del trabajo y contrastado con las fuentes bibliográficas utilizadas.

    \item \textbf{Validación:} La información técnica fue verificada utilizando papers académicos, el material entregado por el curso y juicio propio. La herramienta se utilizó como apoyo de redacción, no como única fuente de información.
\end{enumerate}